\documentclass[a4paper]{arrowhead}

\usepackage[yyyymmdd]{datetime}
\usepackage{etoolbox}
\usepackage[utf8]{inputenc}
\usepackage[T1]{fontenc}
\usepackage{multirow}
\usepackage{longtable}
\usepackage{tikz}
\usetikzlibrary{positioning,shapes.multipart,arrows.meta}

\renewcommand{\dateseparator}{-}
\setlength{\emergencystretch}{3em}
\lstset{language=HTTP,breakatwhitespace=false,columns=fullflexible}

%% Special references
\newcommand{\mref}[2][]{{\textcolor{ArrowheadPurple}{\hyperref[sec:model:#2]{\ifstrempty{#1}{#2}{#1}}}}}
\newcommand{\pref}[1]{{\textcolor{ArrowheadGrey}{\hyperref[sec:model:primitives:#1]{#1}}}}
%%

\begin{document}
\sloppy

%% Arrowhead Document Properties
\ArrowheadTitle{AIGym}
\ArrowheadType{System Design Description}
\ArrowheadTypeShort{SysDD}
\ArrowheadVersion{5.2.2}
\ArrowheadDate{2026-10-30}
\ArrowheadAuthor{Jerker Delsing}
\ArrowheadStatus{RELEASE CANDIDATE}
\ArrowheadContact{jerker.delsing@ltu.se}
\ArrowheadFooter{\href{www.arrowhead.eu}{www.arrowhead.eu}}
\ArrowheadSetup
%%

%% Front Page
\begin{center}
  \vspace*{1cm}
  \huge{\arrowtitle}

  \vspace*{0.2cm}
  \LARGE{\arrowtype}
  \vspace*{1cm}

  \vspace*{\fill}

  % Front Page Image

  \vspace*{1cm}
  \vspace*{\fill}

  % Front Page Abstract
  \begin{abstract}
    This document describes the design and implementation of the AIGym support system of Arrowhead 5.2.2 (\texttt{ah5-AIGym-java-spring}).
  \end{abstract}

  \vspace*{1cm}
\end{center}

\newpage
%%

%% Table of Contents
\tableofcontents
\newpage
%%

\section{Overview}
\label{sec:overview}

This document describes the AIGym system implementation, which realizes the AIGym SysD \cite{aig-sysd}. AIGym is an experimentation system for machine learning in the local cloud: it collects training samples, predicts with k nearest neighbours and consumes the \texttt{sysMonitorData} stream.
In Section \ref{Implementation}, we describe implementation details of the system.

\newpage

\section{Implementation}
\label{Implementation}

This implementation is based on the SysD document AIGym \cite{aig-sysd} and is located in \texttt{ah5-AIGym-java-spring}.

\subsection{Implementation language and tools}
\begin{itemize}
\item Java 21, Spring Boot (web), Eclipse Paho MQTT, RabbitMQ AMQP client.
\item Storage: in memory.
\end{itemize}

\subsection{Independent source and build ownership}
The project folder owns its application source, Maven build, tests, executable JAR, Dockerfile, configuration and Compose definition. No source or image is taken from a sibling support folder or JavaSupportRuntime. Build with Java 21 and \texttt{mvn -B clean verify}; a populated external Maven dependency cache permits offline builds.

AIGym retains its independent Spring Boot parent and application source. It does not use the Arrowhead common runtime or another support-system artifact.

The packaged artifact is \texttt{arrowhead-\allowbreak{}ai-\allowbreak{}gym-5.2.2.jar}. The image builds and tests this folder only; the system has a fixed application main class.

\subsection{Functional properties implementation}

The following components implement the system:

\noindent\begin{longtable}{| p{5.6cm} | p{9.6cm} |} \hline
\rowcolor{gray!33} Component & Responsibility \\ \hline \endhead
\texttt{AiGymController} & HTTP operations. \\ \hline
\texttt{ModelStore} & k-NN store. \\ \hline
\texttt{AiGymRuntime} & MQTT intake and sysMonitor. \\ \hline
\texttt{ArrowheadGateway} & Registration, identity, orchestration, authorization. \\ \hline
\texttt{Arrowhead\allowbreak{}Security\allowbreak{}Filter,\allowbreak{} Crl\allowbreak{}Revocation\allowbreak{}Config} & Caller authentication and CRL. \\ \hline
\end{longtable}

Every request passes \texttt{Arrowhead\allowbreak{}Security\allowbreak{}Filter} (identity token verified with Authentication and matched to the certificate; Consumer\allowbreak{}Authorization \texttt{verify} for the operation) before the controller.

\subsubsection{Data handled and stored}

Models and observations in memory.

The data models of the interfaces are Java records (of \texttt{eu.arrowhead.dto} in the common library, or of the module); the IDD documents name them without the \texttt{DTO} suffix. The records used by this system are: \texttt{Prediction}, \texttt{PredictionRequest}, \texttt{Sample}, \texttt{SampleRequest}.

\subsection{Non functional properties implementation}

\subsubsection{Security}

Spring TLS with mandatory client certificates and \texttt{Crl\allowbreak{}Revocation\allowbreak{}Config}; \texttt{Arrowhead\allowbreak{}Security\allowbreak{}Filter} binds the identity token to the certificate.

\subsubsection{Power management}

None; the system is a server process without power states. The \texttt{sysMonitor} \texttt{power} observation always reports line supply.

\subsubsection{Internal monitoring}

Logs go to stdout. The system publishes its \texttt{sysMonitor} observations every 30 s (see the sysMonitor IDD).

\subsubsection{Configuration}

Configuration is read from environment variables (listed in the SysD); secrets are read from files (\texttt{*\_\allowbreak{}FILE} variables or mounted credential bundles).

\subsection{Deployment}

The local \texttt{compose.yaml} selects only this support system. It joins an operator-selected external cloud network; it does not start core systems or another support system. Copy \texttt{.env.example}, provision this system's own identity and credential files, validate the Compose definition, then run \texttt{docker compose up -d --build} from this folder. The image uses the secure external \texttt{config/application.properties}, mutual TLS, outsourced authentication, certificate/token binding and a mounted CRL. Secrets are read through Spring's file configuration tree. The app runs as UID 10001 with a read-only root filesystem and bounded non-executable temporary/log storage.

Registry, Authentication, DSO, ConsumerAuthorization and required brokers/providers remain network dependencies. Independent deployment does not remove those service contracts. The listener health check is only a liveness check; inspect successful application startup, registration and authenticated service access before acceptance. The local \texttt{DEPLOYMENT.md} supplies the complete procedure. The previously recorded live-cloud deployment is historical baseline evidence; the new packaging does not change a live cloud.

\subsection{Local modifications}
\label{sec:local}

Notes on this implementation in the workspace:
\begin{itemize}
  \item Workspace system (5.2.2 naming).
\end{itemize}

\subsection{Verification}

Independent folder-only Maven builds run the project's JUnit suite. Artifact inspection checks the system main/source and absence of JavaSupportRuntime or other support-system implementations. Images and project-local Compose definitions are tested with private disposable state and generated test identities. Spring startup and authenticated HTTPS requests use isolated mutual-TLS core-service fixtures; missing-token and missing-certificate requests are refused. The distribution verification report records exact results and limitations. Broker, NTP and inter-cloud acceptance remain separate from deployment/startup validation.

\newpage

\section{Services}
\label{sec:services}

The implementation of the services is based on the SD and IDD documents referenced below.

\noindent\begin{longtable}{| p{5.4cm} | p{2.2cm} | p{2.2cm} | p{4.8cm} |} \hline
\rowcolor{gray!33}\textbf{Services produced} & \textbf{SysD ref} & \textbf{SD ref} & \textbf{IDD ref} \\ \hline \endhead
\texttt{aiGym} & \cite{aig-sysd} & \cite{aig-aiGym-sd} & \cite{aig-aiGym-http-idd} \\ \hline
\texttt{sysMonitor} & \cite{aig-sysd} & \cite{aig-sysMonitor-sd} & \cite{aig-sysMonitor-amqp-idd} \\ \hline
&&& \\ \hline
\rowcolor{gray!33}\textbf{Services consumed} & \textbf{SysD ref} & \textbf{SD ref} & \textbf{IDD ref} \\ \hline
systemDiscovery, serviceDiscovery & \cite{core-sr-sysd} & \cite{core-sr-systemDiscovery-sd} & -- \\ \hline
identity & \cite{core-auth-sysd} & \cite{core-auth-identity-sd} & -- \\ \hline
authorization & \cite{core-ca-sysd} & \cite{core-ca-authorization-sd} & -- \\ \hline
service\allowbreak{}Orchestration & \cite{core-dso-sysd} & \cite{core-dso-serviceOrchestration-sd} & -- \\ \hline
sysMonitorData & \cite{smo-sysd} & \cite{smo-sysMonitorData-sd} & -- \\ \hline
sysMonitor (publisher) & AMQP broker (SysMonitor) & -- & -- \\ \hline
\end{longtable}


\bibliographystyle{IEEEtran}
\bibliography{bibliography}

\newpage

\section{Revision History}
\subsection{Amendments}

\noindent\begin{tabularx}{\textwidth}{| p{1cm} | p{3cm} | p{2cm} | X | p{4cm} |} \hline
\rowcolor{gray!33} No. & Date & Version & Subject of Amendments & Author \\ \hline

1 & 2026-10-22 & \arrowversion & Document created from the Arrowhead LaTeX templates for the 5.2.2 Java Spring implementation, including the local extensions of this workspace. & Jerker Delsing \\ \hline
2 & 2026-10-25 & \arrowversion & Self-contained build, private library source, configuration, images, database and isolated deployment verification. & Codex (draft) \\ \hline
3 & 2026-10-30 & \arrowversion & Checked against the primary Java source tree (Eclipse-Arrowhead-java): services, operations, paths and data models unchanged. & Claude Opus 5.5 (draft) \\ \hline
\end{tabularx}

\subsection{Quality Assurance}

\noindent\begin{tabularx}{\textwidth}{| p{1cm} | p{3cm} | p{2cm} | X |} \hline
\rowcolor{gray!33} No. & Date & Version & Approved by \\ \hline

1 & & \arrowversion  &  \\ \hline

\end{tabularx}

\end{document}
